% concatenated from manuscript.tex
\documentclass[11pt]{article}

\usepackage{natbib}
\usepackage{graphicx}%
\usepackage{multirow}%
\usepackage{amsmath,amssymb,amsfonts}%
\usepackage{amsthm}%
\usepackage{mathrsfs}%
\usepackage[title]{appendix}%
\usepackage{xcolor}%
\usepackage{textcomp}%
\usepackage{manyfoot}%
\usepackage{booktabs}%
\usepackage{algorithm}%
\usepackage{algorithmicx}%
\usepackage{algpseudocode}%
\usepackage{listings}%
\usepackage{a4wide}

\usepackage{bbm}
\usepackage{bm}
\usepackage{subcaption}
\usepackage{algorithm}
\usepackage{algpseudocode}

\newcommand{\E}{\mathbb E}
\newcommand{\e}{\mathrm e}
\newcommand{\Var}{\mathrm{Var}}
\newcommand{\F}{\mathcal F}
\newcommand{\D}{\mathrm{d}}

\DeclareMathOperator*{\argmax}{arg\,max}
\DeclareMathOperator*{\argmin}{arg\,min}

\theoremstyle{definition}
\newtheorem{theorem}{Theorem}
\newtheorem{lemma}[theorem]{Lemma}
\newtheorem{corollary}[theorem]{Corollary}
\newtheorem{proposition}[theorem]{Proposition}
\newtheorem{remark}[theorem]{Remark}
\newtheorem{assumption}{Assumption}
\newtheorem{definition}{Definition}
\newtheorem{example}{Example}
\newtheorem{fact}{Fact}
\usepackage{authblk}   % author / affiliation formatting
%opening


\begin{document}
\title{Long-memory Markov chains with power-law intensities}

\author[1]{Kyungsub Lee}

\affil[1]{Department of Statistics, Yeungnam University,
	Gyeongsan, Republic of Korea}

\maketitle

\begin{abstract}
	We introduce a self-exciting point process with power-law intensity dynamics that admits a finite-dimensional Markovian state representation.
	The model is constructed to preserve the local jump and slope update structure of power-law Hawkes processes, while replacing global history dependence with a nonlinear Markov chain governing the intensity dynamics.
	Within a general state-space framework, we establish irreducibility, aperiodicity, and the T-chain property under mild regularity conditions on the inter-arrival time distribution.
	Under an explicit stability condition, we further prove that the latent state process is positive Harris recurrent, ensuring the existence of a unique invariant distribution.
	Simulation results based on the local Whittle estimator show that the proposed Markovian intensity model exhibits long-memory behavior near the boundary of the stability region.
\end{abstract}

\section{Introduction}

This paper introduces a class of Markov chains exhibiting long-memory behavior, driven by intensities that decay according to a power law. 
The proposed process is self-exciting in the sense that each event triggers an increase in the intensity, which subsequently relaxes over time according to a power-law decay. 
By augmenting the states to include the slope of the intensity at the moment of excitation, the resulting dynamics form a fundamentally two-dimensional Markov chain that is nevertheless capable of generating long-range dependence.

Self-exciting intensities with a decaying structure provide a stable framework for modeling event-time dynamics in a wide range of natural and social systems.
The Hawkes process \citep{Hawkes1, Hawkes2} is the canonical example, in which each event increases the intensity and the excitation decays over time according to a prescribed kernel.
Exponential decay leads to rapid attenuation and short-range dependence, whereas power-law decay induces slow relaxation and persistent temporal dependence \citep{hardiman2013critical, bacry2016estimation, zhang2016modeling, nystrom2022hawkes}.

Despite their empirical appeal, Hawkes processes with power-law kernels lack a finite-dimensional Markovian representation.
Unlike exponential-kernel Hawkes processes, which admit a Markov formulation via intensity state, power-law kernels induce dependence on the entire event history.
This inherent non-Markovian structure complicates theoretical analysis and practical tasks such as likelihood-based inference, simulation, and control.
As a result, existing work often relies on approximations or alternative constructions, including kernel truncation, sums of exponentials, and nonparametric methods.

Several approaches have been proposed to mitigate the non-Markovian nature of Hawkes processes with power-law kernels. 
\cite{bacry2011nonparametric} introduce a nonparametric estimation method based on second-order statistics, which avoids parametric kernel specification but retains full historical dependence. 
\cite{zhang2022efficient} develop computationally efficient inference procedures using basis expansions of the kernel, effectively approximating infinite memory with a finite-dimensional representation. 
\cite{Zhuang2020NonparametricFramework} proposes a general nonparametric framework for Hawkes kernel estimation that emphasizes flexibility over Markovian structure. 
\cite{bonnet2025nonparametric} further refine nonparametric Hawkes modeling through regularization and improved estimation techniques, while remaining non-Markovian. 
\cite{khabou2025markov} explicitly target Markovian approximations by introducing finite-dimensional latent-state representations for general Hawkes kernels.

Long memory in inter-event durations has also been modeled within the autoregressive conditional duration (ACD) framework of \cite{EngleRussell1998}. 
Drawing on the ARFIMA approach of \cite{granger1980longmemory}, a number of extensions have been proposed; we mention only a few representative examples. 
For instance, \cite{deo2010longmemory} provide semiparametric evidence of long memory in intertrade durations of NYSE stocks and propose a latent-variable long-memory stochastic duration model, 
while \cite{feng2015forecasting} develop a fractionally integrated Log-ACD model with a semiparametric scale function for forecasting trading volumes and related quantities. 
These models capture long-range dependence in event timing through fractional integration, but do not admit a finite-dimensional Markovian representation.

This paper introduces a self-exciting point process that combines long-memory behavior with finite-dimensional Markovian dynamics.
The model retains the local update mechanism of power-law Hawkes processes at event times, capturing both intensity jumps and slope updates, while replacing global history dependence with a 2-dimensional latent state.
Between events, the conditional intensity decays according to a power law driven by the latent state, which is updated at each arrival in a manner consistent with power-law excitation.
The resulting process is Markovian and reproduces key features of power-law self-excitation without requiring the full event history.

Related work has explored Hawkes-type models augmented with additional state variables to encode memory or inhibition effects.
\cite{CostaMaillardMuraro2024} analyze a discrete-time Hawkes process with inhibition and finite-length memory, providing a near-complete stability characterization.
Related Markovian constructions include the renewal-based
stability analysis of signed Hawkes processes in
\cite{costa2020renewal} and the ephemerally self-exciting
process of \cite{daw2022ephemerally}, which confines each
event's excitation to a random duration; our model instead
encodes power-law decay in a fixed two-dimensional latent state.

Variable-length memory formulations are studied in \cite{hodara2017hawkes} and more recently in \cite{quayle2025hawkesprocessesvariablelength}, where the conditional intensity depends on a dynamically selected portion of the past.
\cite{Schmutz2022MeanField} studies a mean‑field Hawkes process with age and leaky memory, proving propagation of chaos in the large‑population limit.
These approaches retain explicit dependence on past events, leading to higher-order or variable-dimensional state representations.
In contrast, our model replaces explicit history dependence with a fixed-dimensional latent Markov state, preserving power-law self-excitation while yielding a tractable continuous-time formulation.
 
The proposed model can be expressed within a general inverse-integrated intensity framework, in which inter-arrival times are generated by applying an inverse transform to an independent sequence through the integrated conditional intensity as in~\cite{lee2025self,lee2026forecasting}.
This formulation enables exact simulation, likelihood-based inference, and tractable analysis.

The process further admits a rigorous stability and ergodicity analysis within a state-space Markov chain framework. 
Following the general theory of Markov chains developed in \cite{meyn2009markov}, 
we formulate the dynamics as a nonlinear stochastic system driven by exogenous inter-arrival times and analyze the associated control model.
The embedded Markov chain is shown to be $\psi$-irreducible, aperiodic,
and positive Harris recurrent under the stability condition, 
implying the existence of a unique invariant distribution and long-run stability.

Our approach aligns with a broader line of work that leverages this framework to obtain invariant measure existence/uniqueness and quantitative convergence properties for Markovian state-space models; see, e.g., \citet{MayerhoferStelzerVestweber2020AffineCones,ShaoWangWu2023HybridSwitching}. 
In the Hawkes literature, similar recurrence-and-stationarity arguments for Markovian representations of the intensity have also been used as a key technical input for asymptotic analysis \citep{PalmowskiPojerThonhauser2025RuinHawkes} and for ergodic characterization under general excitation kernels \citep{KwanChenDunsmuir2025ErgodicGeneralKernelHawkes}.

The remainder of the paper is organized as follows.
Section~\ref{Sect:motive} reviews the motivation based on power-law Hawkes processes.
Section~\ref{Sect:model} introduces the proposed Markovian power-law intensity model and its structural properties.
Section~\ref{Sect:MC} discusses the Markov properties of the model.
Section~\ref{Sect:simul} presents simulation results and long-memory analysis.
Section~\ref{Sect:conc} concludes.

\section{Motivation}\label{Sect:motive}

A primary motivation is the Hawkes process with a power-law kernel, which exhibits long-memory behavior.
Let $\{T_n\}_{n\ge1}$ denote the event times.
Consider a Hawkes process with conditional intensity
\[
\lambda(t)=\mu+\sum_{T_i<t} g(t-T_i),
\qquad 
g(u)=\frac{\alpha}{(u+\gamma)^p}, \quad u>0,
\]
where $\mu>0$, $\alpha>0$, $\gamma>0$, and $p>0$.

For $p>1$, the kernel is integrable,
\[
\int_0^\infty g(u)\,\D u<\infty,
\]
and the process is stationary when the branching ratio
$\int_0^\infty g(u)\,\D u<1$, i.e.,
\[
\frac{\alpha}{p-1}\gamma^{1-p}<1.
\]
When $1<p<2$, the kernel decays polynomially,
\[
g(u)\sim \alpha u^{-p}, \qquad u\to\infty,
\]
leading to slow temporal decay.
In this regime, the autocovariance of the counting process satisfies
\[
\mathrm{Cov}\!\left(N(t+h)-N(t),\,N(0,1]\right)
\sim C h^{-(p-1)}, \qquad h\to\infty,
\]
for some $C>0$.

At an event time $T_n$, the conditional intensity satisfies
\begin{align}
	\lambda(T_n^+)-\lambda(T_n^-)
	&= g(0)=\frac{\alpha}{\gamma^p}, \label{Eq:update1}\\
	\lambda'(T_n^+)-\lambda'(T_n^-)
	&= g'(0)=-\frac{\alpha p}{\gamma^{p+1}}. \label{Eq:update2}
\end{align}
Thus, each arrival induces a constant upward jump in the intensity and a constant downward jump in its slope.
Between events, both quantities decay according to the power-law kernel.

Unlike the exponential-kernel Hawkes process, the power-law Hawkes process does not admit a finite-dimensional Markovian representation.
In the exponential case, superposed contributions share a common decay rate, yielding a closed Markovian state.
For power-law kernels, the decay depends on elapsed time since each event, so the intensity depends on the entire history.
Our objective is to construct a Markov process that preserves the local update structure in
Eqs.~\eqref{Eq:update1}--\eqref{Eq:update2}.



\section{Model}\label{Sect:model}	

Let \(\{T_n\}_{n \ge 1}\) denote the sequence of event arrival times.
For \(T_n < t' \le T_{n+1}\), define the local time \(t = t' - T_n\).
The conditional intensity is given by
\begin{equation}
	\lambda(t') = \lambda_{n}(t) = \mu + \frac{a_n}{(t  + c_n)^p},
	\label{eq:intensity}
\end{equation}
where \(\mu > 0\) is the baseline intensity, \(p > 0\) is the fixed decay exponent, and \(a_n, c_n > 0\) are stochastic processes determined at  $T_n$.

At the \(n\)-th event time \(T_n\), given the inter-arrival time \(\tau_n = T_n - T_{n-1}\), the processes \(a_n\) and \(c_n\) are updated to satisfy the following  conditions
\begin{align*}
	& \lambda_{n}(0) = \lambda_{n-1}(\tau_{n}) + \frac{\alpha}{\gamma^p} \\
	&\frac{\D \lambda_{n}}{\D t} (0) = \frac{\D \lambda_{n-1}}{\D t}(\tau_{n})  + \frac{\D}{\D t} \left. \frac{\alpha}{(\gamma + t)^p} \right|_{t = \tau_n} = \frac{\D \lambda_{n-1}}{\D t}(\tau_{n})   - \frac{p \alpha}{\gamma^{p+1}}
\end{align*}
where \(\alpha\) and \(\gamma\) are model parameters.
These equations ensure 
that our model follows the update rule in Eqs.~\eqref{Eq:update1}~and~\eqref{Eq:update2}
as in the Hawkes process with power-law kernel.
Consequently,
\begin{align}
	&\frac{a_n}{c_n^p} = \frac{a_{n-1}}{(\tau_n + c_{n-1})^p} + \frac{\alpha}{\gamma^p} \label{eq:update_level}\\
	&\frac{a_n}{c_n^{p+1}} = \frac{a_{n-1}}{(\tau_n + c_{n-1})^{p+1}} + \frac{\alpha}{\gamma^{p+1}}. \label{eq:update_slope}
\end{align}
Solving the equation yields
\begin{align*}
	c_n &= \frac{a_{n-1}\gamma^{p+1}\,(\tau_n+c_{n-1}) \;+\; \alpha\gamma\,(\tau_n+c_{n-1})^{p+1}}
	{a_{n-1}\gamma^{p+1} \;+\; \alpha\,(\tau_n+c_{n-1})^{p+1}}, \\
	a_n &= c_n^p  \left( \frac{a_{n-1}}{(\tau_n + c_{n-1})^p} + \frac{\alpha}{\gamma^p} \right).
\end{align*}

Meanwhile, an alternative parameterization is convenient for analysis.
Let
$$ X_n = \frac{a_n}{c_n^p}, \quad \xi = \frac{\alpha}{\gamma^p}.$$
Then the conditional intensity function can be written as
\begin{equation}
\lambda_n(t; X_n, c_n) = \mu + X_n \left(1 + \frac{t}{c_n} \right)^{-p}. \label{eq:intensity2}
\end{equation}
Under this parameterization, the update rules become
\begin{align*}
	&X_n = X_{n-1} \left(1  + \frac{\tau_n}{c_{n-1}} \right)^{-p} + \xi \\
	&\frac{X_n}{c_n} = \frac{X_{n-1}}{c_{n-1}} \left(1  + \frac{\tau_n}{c_{n-1}} \right)^{-p-1} + \frac{\xi}{\gamma}.
\end{align*}
Let
$$r_n = 1 + \frac{\tau_n}{c_{n-1}},$$
then $r_n^{-p}$ and  $r_n^{-p-1}$ act as decay factors, satisfying
$0 < r_n^{-p}, r_n^{-p-1} < 1$.
With this notation, $c_n$ can be expressed as
$$
c_n
\;=\;
\frac{X_{n-1}\,r_n^{-p} + \xi}{\dfrac{X_{n-1}}{c_{n-1}}r_n^{-(p+1)} + \frac{\xi}{\gamma} }.
$$

We now formalize this construction.
We begin with a general modeling framework, adapted from \cite{lee2025self}.

\begin{definition}\label{Def:general_vector}
	Let $\{ \varepsilon_n \}_{n \in \mathbb{N}}$ be a sequence of independent and identically distributed random variables. 
	Suppose there exists a stochastic process $\{ \mathbf{X}_n \}$ taking values in $\mathbb{R}^d$ for some $d \ge 1$, which governs the distribution of inter-arrival times $\{ \tau_n \}$.
	Let \(\Phi(t,\mathbf{x})\) be a function defined for \(t \ge 0\) and \(\mathbf{x} \in \mathbb{R}^d\), and assume that \(\Phi(\cdot,\mathbf{x})\) is strictly increasing in \(t\) for each fixed \(\mathbf{x}\).
	The inter-arrival time \(\tau_n\) is related to \(\varepsilon_n\) by
	\begin{equation}
		\varepsilon_n = \Phi(\tau_n, \mathbf{X}_{n-1}) 
		\qquad \text{equivalently,} \qquad
		\tau_n = \Phi^{-1}(\varepsilon_n, \mathbf{X}_{n-1}) ,
		\label{Eq:tau_vector}
	\end{equation}
	where the inverse is taken with respect to $t$, holding $\mathbf{x}$ fixed.
	The process $\mathbf{X}_n$ evolves according to an update function $\Psi$ such that
	\begin{equation}
		\mathbf{X}_n = \Psi(\mathbf{X}_{n-1}, \tau_n). \label{Eq:Psi0}
	\end{equation}
	The associated point process \(N\) on \(\mathbb{R}^+\) is defined by
	\[
	N = \sum_{n \in \mathbb N} \delta_{T_n}, 
	\qquad  \qquad 
	T_n = \sum_{i=1}^n \tau_i
	\]
	where \(\delta_t\) denotes the Dirac measure at \(t\).
\end{definition}

We now specialize this framework to the model of interest, 
in which the latent rate function is self-exciting and decays 
according to a power law between successive events, 
with the decay governed by the state variables $(X_n, c_n)$.

\begin{definition}\label{Def:Model}
	The latent processes $\mathbf{X}_n = \{(X_n, c_n)\}$ evolve according to 
	\begin{equation}
		\mathbf{X}_n = \Psi(\mathbf{X}_{n-1}, \tau_n) = (\Psi_x(x, c, t), \Psi_c(x, c, t)).
		\label{Eq:Psi}
	\end{equation}
	where
	\begin{align}
		\Psi_x(x, c, t) &= x \left(1 + \frac{t}{c}\right)^{-p} + \xi,  \label{Eq:Psi_x}\\
		\Psi_c(x, c, t) &= \frac{\Psi_x(x, c, t)}{\,\frac{x}{c} \left(1 + \tfrac{t}{c}\right)^{-p-1} + \tfrac{\xi}{\gamma}}, \label{Eq:Psi_c}
	\end{align}
	with parameters $\mu> 0, \gamma> 0, \xi > 0$ and $p>1$.
	Equivalently, the update equations can be written as
	\begin{align*}
		&X_n = X_{n-1} \left(1  + \frac{\tau_n}{c_{n-1}} \right)^{-p} + \xi \\
		&\frac{X_n}{c_n} =  \frac{X_{n-1}}{c_{n-1}} \left(1  + \frac{\tau_n}{c_{n-1}} \right)^{-p-1} + \frac{\xi}{\gamma}.
	\end{align*}
	Define $\Phi(t, x, c)$ as
	\begin{align}
		\Phi(t, x, c) & = \int_0^t  \lambda_{n-1}(s; x, c) \D s= \int_0^t \left( \mu + \frac{x c^p}{(c+s)^p} \right)\, \D s \label{Eq:Phi1}\\
		& = \mu t + \dfrac{x c}{p-1}\left( 1 - \bigg(1+\dfrac{t}{c}\bigg)^{-p+1}  \right). \label{Eq:Phi}
	\end{align}
	The inter-arrival time $\tau_n$ is generated via the inverse transform
	\[
	\tau_n = \Phi^{-1}(\varepsilon_n,  X_{n-1}, c_{n-1}),
	\]
	where the inverse is taken with respect to $t$ for fixed $x$ and $c$,
	and $\{\varepsilon_n\}$ is an i.i.d. sequence drawn from $F_\varepsilon$
	satisfying $F_\varepsilon(0) = 0$ (so that $\varepsilon_n \ge 0$ a.s.), $\mu_{\varepsilon} = \mathbb{E}[\varepsilon] < \infty$,
	and $\mathbb{E}[\varepsilon^2] < \infty$.
\end{definition}

Here, $\lambda$ in Eq.~\eqref{Eq:Phi1} plays the role 
of the conditional intensity when $F_\varepsilon$ is the unit exponential 
distribution; in the general case, it is not a conditional intensity 
in the strict sense, but serves as the state-driven rate function 
governing the shape of the inter-arrival time distribution.

Let $\E_{x,c}[\,\cdot\,]$ denote the conditional expectation given $X_{n-1} =x$ and $c_{n-1} = c$.
The model relation
$\Phi(\tau_n,x,c)=\varepsilon_n$ implies
\[
F_{\tau}(s)=\mathbb{P}(\tau_n\le s)
=F_{\varepsilon}(\Phi(s,x,c)),
\qquad
f_{\tau}(s)=f_{\varepsilon}(\Phi(s,x,c))\,\frac{\partial \Phi(s,x,c)}{\partial s}.
\]
Hence,
\[
\mathbb{E}_{x,c}[\tau_n]
=\int_0^{\infty} s\, f_{\varepsilon}(\Phi(s,x,c))\, \frac{\partial \Phi(s,x,c)}{\partial s}\, \D s.
\]
Let $\bar F$ be a survival function corresponding to a generic distribution $\F$. 
By integration by parts,
$$
\begin{aligned}
	\mathbb{E}_{x,c}[\tau_n] &= \left[ -s \bar{F}_{\tau}(s) \right]_{0}^{\infty} + \int_{0}^{\infty} \bar{F}_{\tau}(s) \D s \\
	&= \left[ -s \mathbb P(\tau_n > s) \right]_{0}^{\infty} + \int_{0}^{\infty} \mathbb P(\tau_n > s) \D s \\
	& =\int_0^{\infty}\mathbb{P}(\tau_n>s)\, \D s,
\end{aligned}
$$
where the boundary term vanishes  since 
$\mathbb{E}_{x,c}[\tau_n] < \infty$ under the moment conditions 
$\mu_\varepsilon < \infty$ imposed on $F_\varepsilon$.
Using
$\mathbb{P}(\tau_n>s)=\bar F_{\varepsilon}(\Phi(s,x,c))$, we obtain
\begin{equation}
	\mathbb{E}_{x,c}[\tau_n]
	=\int_0^{\infty}\bar F_{\varepsilon}(\Phi(t,x,c))\,\D t \label{Eq:CE_tau}
\end{equation}
and
\begin{equation}
\mathbb{E}_{x,c}[\tau_n^k] = \int_0^\infty k t^{k-1} \bar{F}_\varepsilon(\Phi(t, x, c)) \D t. \label{Eq:CM_tau}
\end{equation}

Using $f_\tau(s) = f_\varepsilon(\Phi(s,x,c))\,\lambda_n(s;x,c)$, 
the log-likelihood based on $\{\tau_n\}_{n=1}^N$ is
\begin{equation}
	\log L
	=
	\sum_{n=1}^N
	\left[
	\log f_\varepsilon\!\left(\Phi(\tau_n, X_{n-1}, c_{n-1})\right)
	+
	\log \lambda_n(\tau_n; X_{n-1}, c_{n-1})
	\right].
\end{equation}



\subsection{Properties of $X$ and $c$}


\begin{proposition}
	The one-step update $c_n$ is bounded.
\end{proposition}
\begin{proof}
	Fix $c_{n-1} = c$ and $X_{n-1} = x$. 
	Recall that the update formula for $c_n$ can be written as a function of the growth ratio $r = 1 + \frac{\tau}{c}$ in the form
	$$c_n(r) = \frac{x r^{-p} + \xi}{\frac{x}{c} r^{-(p+1)} + \frac{\xi}{\gamma}}, \quad r \in [1, \infty)$$
	where $p > 1$ and $\xi, \gamma, c, x > 0$.
	Define the numerator and denominator as
	\[
	N(r) = x r^{-p} + \xi,
	\qquad
	D(r) = \frac{x}{c} r^{-(p+1)} + \frac{\xi}{\gamma}.
	\]
	Then
	$$c_n'(r) = \frac{N'(r)D(r) - N(r)D'(r)}{[D(r)]^2}.$$
	A direct calculation shows that the numerator equals
	\begin{equation}
	\frac{x^2}{c}r^{-2p-2} - \frac{p \xi x}{\gamma}r^{-(p+1)} + \frac{(p+1)\xi x}{c}r^{-(p+2)}.\label{Eq:eq1}
	\end{equation}
	Thus, the condition $c_n'(r) = 0$ is equivalent to finding $r$ such that
	$$H(r) = - c p r + \gamma(p+1) + \frac{\gamma x}{\xi} r^{-p} = 0,$$
	which is obtained by multiplying Eq.~\eqref{Eq:eq1} by $\frac{\gamma c}{\xi x} r^{p+2}$.
	Let $r^*$ be the solution of $H(r) =0$.
	Note that
	$$ H'(r) = -\frac{p \gamma x}{\xi} r^{-(p+1)} - cp < 0, \qquad r \geq 1,$$
	which implies $H(r)$ is strictly decreasing on \([1,\infty)\).
	Thus, either $H(r)$ is all negative for all $r \geq 1$, or $H(r) > 0$ for $r < r^*$ and $H(r) < 0$ for $r > r^*$.
	Moreover, the sign of \(H(r)\) coincides with that of \(c_n'(r)\).
	Consequently, either \(c_n(r)\) is strictly decreasing on \([1,\infty)\), or there exists a unique maximizer \(r^* \ge 1\) such that \(c_n(r)\) is increasing for \(r<r^*\) and decreasing for \(r>r^*\).
	In either case, the supremum of \(c_n\) is attained at
	\begin{equation}
	\sup c_{n} = \max \{ c_n(1), c_n(r^*) \}. \label{Eq:max_cn}
	\end{equation}
	This completes the proof.
\end{proof}


\begin{lemma}\label{lem:mean-tau-bound}
	Let $I(K) = \int_{K}^{\infty} \bar{F}_{\varepsilon}(u) \D u$. 
	For all $x,c>0$, the inter-arrival time $\tau_n$ satisfies
	\begin{equation}
		\mathbb{E}_{x,c}[\tau_n] \le \frac{\mu_{\varepsilon}}{\mu + k_p x} + \dfrac{I\left((\mu + k_p x)c\right)}{\mu}, \quad\quad k_p = \frac{1-2^{\,1-p}}{p-1}.
	\end{equation}
	For any $x_* > 0$, the conditional expectation $\mathbb{E}_{x,c}[\tau_n]$ satisfies:
	\[
	\mathbb{E}_{x,c}[\tau_n] \le 
	\begin{cases} 
		\dfrac{\mu_{\varepsilon}}{\mu + k_p x} + \dfrac{I((\mu + k_p x)c)}{\mu}, & x \ge x_*, \\[3ex]
		\dfrac{\mu_{\varepsilon}}{\mu} + \dfrac{I(\mu c)}{\mu}, & 0 \le x < x_*.
	\end{cases}
	\]
\end{lemma}


\begin{proof}
	Recall that
	\[
	\mathbb E_{x,c}[\tau_n]
	=
	\int_0^\infty \bar{F}_{\varepsilon}(\Phi(t;x,c))\, \D t
	\]
	where
	\[
	\Phi(t;x,c)
	=
	\mu t
	+
	\frac{x c}{p-1}
	\left[
	1-\left(1+\frac{t}{c}\right)^{1-p}
	\right].
	\]
	Since \(\Phi(\cdot\,;x,c)\) is concave and satisfies \(\Phi(0;x,c)=0\), it lies above the secant line connecting \((0,0)\) and \((c,\Phi(c;x,c))\) on the interval \([0,c]\).
	Evaluating at \(t=c\), we obtain
	\[
	\Phi(c;x,c) = \mu c + \frac{xc}{p-1}(1 - 2^{1-p}) = c(\mu + k_p x).
	\]
	Therefore, for \(0 \le t \le c\), 
	\begin{equation}
	\Phi(t;x,c) \ge \frac{\Phi(c;x,c)}{c} t = (\mu + k_p x)t. \label{Eq:Phi_B1}
	\end{equation}
	Since $\bar{F}_{\varepsilon}$ is non-increasing, we have
	\begin{equation}
		\int_0^c \bar{F}_{\varepsilon}(\Phi(t;x,c))\, \D t \le
		\int_0^c \bar{F}_{\varepsilon}((\mu+ k_p x)t)\, \D t \le \int_0^\infty \bar{F}_{\varepsilon}((\mu+ k_p x)t)\, \D t = \frac{\mu_{\varepsilon}}{\mu + k_p x}. \label{Eq:bound1_gen}
	\end{equation}
	For $t\ge c$, the intensity is bounded below by $\mu$, which implies
	\begin{equation}
	\Phi(t;x,c)
	=
	\Phi(c;x,c)+\int_c^t \lambda_{n-1}(s;x,c)\,\D s
	\ge
	\Phi(c;x,c)+\mu(t-c) = c(\mu + k_p x) + \mu(t-c). \label{Eq:Phi_B2}
	\end{equation}
	Consequently, using the substitution $u = \Phi(c) + \mu(t-c)$,
	\begin{equation}
		\int_c^\infty \bar{F}_{\varepsilon}(\Phi(t;x,c))\,\D t
		\le \int_c^\infty \bar{F}_{\varepsilon}(\Phi(c) + \mu(t-c))\,\D t = \frac{1}{\mu} \int_{\Phi(c)}^\infty \bar{F}_{\varepsilon}(u) du = \frac{I(c(\mu+k_p x))}{\mu}. \label{Eq:bound2_gen}
	\end{equation}
	Alternatively, for \(t \ge c\), since $\Phi(t) \ge \mu t$,
	\begin{equation}
		\int_c^\infty \bar{F}_{\varepsilon}(\Phi(t;x,c))\,\D t
		\le
		\int_c^\infty \bar{F}_{\varepsilon}(\mu t)\,\D t
		=
		\frac{I(\mu c)}{\mu}. \label{Eq:bound3_gen}
	\end{equation}
	For any fixed $x_*>0$, if $0\le x<x_*$, 
	combining the two bounds \eqref{Eq:bound1_gen} and \eqref{Eq:bound3_gen} 
	yields
	\[
	\mathbb E_{x,c}[\tau_n]
	\le \frac{\mu_{\varepsilon}}{\mu + k_p x} + \frac{I(\mu c)}{\mu} \le
	\frac{\mu_{\varepsilon}}{\mu} + \frac{I(\mu c)}{\mu},
	\]
	If $x > x_*$, combining the two bounds \eqref{Eq:bound1_gen} and \eqref{Eq:bound2_gen} 
	yields
	\[
	\mathbb E_{x,c}[\tau_n]
	\le
	\frac{\mu_{\varepsilon}}{\mu+ k_p x} + \frac{I(c(\mu+k_p x))}{\mu}.
	\]
\end{proof}

\begin{proposition}[Negative drift of $c_n$]\label{Prop:drift_c}
	There exists $x_* > 0$ such that the following statements hold.
	\begin{enumerate}
		\item For $0 < x < x_*$, if 
		\begin{equation}
			c > C_1 
			:= \frac{2\mu \gamma + 3 \mu_{\varepsilon} + \sqrt{(2\mu \gamma + 3 \mu_{\varepsilon})^2 
					+ \frac{24  \mu \mu_{\varepsilon} \gamma x_*}{\xi}}}{4\mu}, \label{Eq:C1}
		\end{equation}
		then $\mathbb{E}_{x,c}[c_n] - c < 0$.
		
		\item For $x \ge x_*$, if 
		\begin{equation}
			c > C_2 
			:= \frac{\gamma + \sqrt{\gamma^2 + \frac{3 \mu_{\varepsilon} \gamma}{k_p \xi}}}{2}, \label{Eq:C2}
		\end{equation}
		then $\mathbb{E}_{x,c}[c_n] - c < 0$.
	\end{enumerate}
\end{proposition}


\begin{proof}
	From the update formula and using $X_{n-1} = x$ and $c_{n-1} = c$, we have
	\[
	c_n - c
	= \frac{x\,r_n^{-p} + \xi}
	{D} - c = 
	\frac{\frac{x \tau_n}{c} r_n^{-(p+1)} + \xi\left(1-c/\gamma\right)}{D}, \quad\quad D=\frac{x}{c} r_n^{-(p+1)}+\frac{\xi}{\gamma}
	\]
	where $r_n = 1 + \tau_n/c$.
	First, note that 
	\[
	\frac{\frac{x\tau_n}{c} r_n^{-(p+1)}}{D}
	\le
	\frac{\frac{x\tau_n}{c} r_n^{-(p+1)}} {\frac{x}{c} r_n^{-(p+1)}}
	= \tau_n.
	\]	
	Thus,
	\begin{equation}
		c_n - c \leq \tau_n + \frac{ \xi(1-c/\gamma)}{D} \label{Eq:c_tau}
	\end{equation}
	which intuitively suggests that if $c$ is large enough then the expectation of the right-hand side is negative.
	Since 
	$$ D=\frac{x}{c} r_n^{-(p+1)}+\frac{\xi}{\gamma} \leq \frac{x}{c} +\frac{\xi}{\gamma}$$
	and $1  - c/\gamma$ is negative when $c > \gamma$, 
	we have
	$$\frac{\xi(1 - c/\gamma)}{D} \le \frac{\xi(1 - c/\gamma)}{\frac{x}{c} + \frac{\xi}{\gamma}} =  \frac{c\xi(\gamma - c)}{\gamma x + c \xi}$$
	since $r_n \ge 1$.
	
	Taking the conditional expectation of Eq.~\eqref{Eq:c_tau} and applying Lemma~\ref{lem:mean-tau-bound}, for any $x_*>0$, we have
	\begin{equation}\label{Eq:U_bound}
		\mathbb{E}_{x,c}[\,c_n - c\,] \leq  
		\begin{cases}
			U_1(x,c)
			= \dfrac{\mu_{\varepsilon}}{\mu + k_p x}
			+ \dfrac{I(c(\mu + k_p x))}{\mu}
			+ \dfrac{c\xi(\gamma - c)}{\gamma x + c \xi},
			& x \ge x_*, \\[2ex]
			U_2(x,c)
			= \dfrac{\mu_{\varepsilon}}{\mu}
			+ \dfrac{I(\mu c)}{\mu}
			+ \dfrac{c\xi(\gamma - c)}{\gamma x + c \xi},
			& 0 \le x < x_* .
		\end{cases} 
	\end{equation}
	
	For the region $x \ge x_*$, we analyze the tail term $I(c(\mu + k_p x))$. Since $\mathbb{E}[\varepsilon^2] < \infty$ implies $I(K) = o(K^{-1})$, the term $I(c(\mu + k_p x))$ decreases faster than $(\mu + k_p x)^{-1}$. We define $C_3(x_*)$ as the smallest $c \ge \gamma$ such that 
	\begin{equation}
		\frac{I(c(\mu + k_p x_*))}{\mu} \le \frac{\mu_{\varepsilon}}{2(\mu + k_p x_*)}.  \label{Eq:I_star}
	\end{equation}
	Because $I(K)$ is monotonically decreasing, for $c > C_3$ and $x \ge x_*$, we have $I(c(\mu+k_px)) \le I(C_3(\mu+k_px_*))$, ensuring that
	\begin{equation}
		\frac{I(c(\mu + k_p x))}{\mu} \le \frac{\mu_{\varepsilon}}{2(\mu + k_p x)}. \label{Eq:I}
	\end{equation}
	holds throughout the region. 
	Substituting Eq.~\eqref{Eq:I} into $U_1(x, c)$ yields
	\begin{equation}
		U_1(x,c) \le  \frac{3 \mu_{\varepsilon}}{2(\mu + k_p x)} + \frac{c\xi(\gamma - c)}{\gamma x + c \xi} \le \frac{3\mu_{\varepsilon}}{2 k_p x} + \frac{c\xi(\gamma - c)}{\gamma x + c \xi}. \label{Eq:U1_bound}
	\end{equation}
	Furthermore, as $x_*$ increases, the right-hand side of Eq.~\eqref{Eq:I_star} decreases at rate $O(x_*^{-1})$, but the argument of $I$ in the left-hand side of Eq.~\eqref{Eq:I_star} increases, 
	making $C_3(x_*)$ a non-increasing function that can be made arbitrarily close to $\gamma$ for sufficiently large $x_*$.
	
	Furthermore, for
	\[
	c > C_4 := \gamma + \frac{3 \mu_{\varepsilon}}{2 k_p} \left( \frac{1}{x_*} + \frac{1}{\xi} \right),
	\]
	we have the following sequence of inequalities:
	\[
	2 k_p x_* (c - \gamma) > 3 \mu_{\varepsilon} + \frac{3 \mu_{\varepsilon} x_*}{\xi} > 3 \mu_{\varepsilon} + \frac{3 \mu_{\varepsilon} \gamma x_*}{\xi c},
	\]
	where we utilize $c > \gamma$ for the second inequality. This leads to
	\[
	3 \mu_{\varepsilon} c \xi < (2 k_p \xi c (c - \gamma) - 3 \mu_{\varepsilon} \gamma) x_* \le (2 k_p \xi c (c - \gamma) - 3 \mu_{\varepsilon} \gamma) x,
	\]
	since the quadratic term $2 k_p \xi c (c - \gamma) - 3 \mu_{\varepsilon} \gamma$ is positive for all $c > C_2$ in Eq.~\eqref{Eq:C2}. 
	Rearranging these terms, we obtain
	\[
	3 \mu_{\varepsilon} (\gamma x + c \xi) < 2 k_p \xi c (c - \gamma) x \implies \frac{3 \mu_{\varepsilon}}{2 k_p x} < \frac{c \xi (c - \gamma)}{\gamma x + c \xi},
	\]
	and by Eq.~\eqref{Eq:U1_bound}, we have $U_1(x, c) < 0$. 
	Since $\lim_{x_* \to \infty} C_3(x_*) = \gamma$ and $\lim_{x_* \to \infty} C_4(x_*) = \gamma + \frac{3\mu_\varepsilon}{2k_p\xi}$, while $C_2$ remains a constant greater than these limits for large $x_*$, the condition $c > C_2$ is sufficient.
	
	For the region $0 \le x < x_*$, 
	let $C_5$ be the infimum of $c$ such that $I(\mu c) \le \mu_{\varepsilon}/2$. Then, for $c > C_5$,
	\[
	U_2(x,c) \le \frac{3\mu_{\varepsilon}}{2\mu} + \frac{c\gamma\xi(1 - c/\gamma)}{\gamma x_* + c \xi}.
	\]
	Solving the quadratic inequality 
	$$\frac{3\mu_{\varepsilon}}{2\mu} (\gamma x_* + c \xi) + c \gamma \xi (1 - c/\gamma) \leq 0$$ for $c$ yields the threshold $C_1$ in Eq.~\eqref{Eq:C1}. 
	Since $C_1$ grows with $x_*$ while $C_5$ is independent of $x_*$, we can always choose a sufficiently large $x_*$ such that $C_1 > C_5$. 
	Under this choice, for all $c > C_1$, the condition $c > C_5$ is automatically satisfied, ensuring that the conditional drift is strictly negative uniformly over $0 \le x < x_*$.
\end{proof}



\begin{proposition}[Negative drift of $X_n$]\label{Pro:ndX}
	Suppose that the model parameters satisfy 
	\begin{equation}
		\xi < \frac{p}{2 \gamma} \label{Eq:condition1}
	\end{equation}
	If $ c < 2 \gamma$, or equivalently $ c < p/\xi$,
	then there exists $x^* > 0$ such that for all $x > x^*$, the conditional drift of $X_n$ is negative, i.e.,
	$$\mathbb{E}_{x, c}[X_n] - x < 0.$$
\end{proposition}

\begin{proof}
	By the second-order Taylor expansion, 
	$$(1+u)^{-p} \le 1 - pu + \frac{p(p+1)}{2}u^2,$$
	for all $u\ge0$ and substituting $u = \tau_n/c$, the change in $X$ is bounded by 
	\begin{equation}\label{eq:drift_base_general}
		\mathbb{E}_{x,c}[X_n - x] \le \xi - \frac{px}{c} \mathbb{E}_{x,c}[\tau_n] + \frac{xp(p+1)}{2c^2} \mathbb{E}_{x,c}[\tau_n^2].
	\end{equation}
	To analyze the moments of $\tau_n$ with $F_\varepsilon$, we use the relations Eqs.~\eqref{Eq:CE_tau}~and~\eqref{Eq:CM_tau}. 
	For the first moment, the concavity of $\Phi$ implies $\Phi(t) \le (\mu + x)t$, thus
	\begin{equation}
		\mathbb{E}_{x,c}[\tau_n] = \int_0^\infty \bar{F}_\varepsilon(\Phi(t,x,c)) \D t \ge \int_0^\infty \bar{F}_\varepsilon((\mu + x)t) \D t = \frac{\mu_{\varepsilon}}{\mu + x}.
	\end{equation}
	Substituting this into the second term of \eqref{eq:drift_base_general} gives $$-\frac{px}{c} \mathbb{E}_{x,c}[\tau_n] \le -\frac{p}{c} \left(\frac{\mu_{\varepsilon} x}{\mu + x}\right).$$ 
	
	For the second moment, applying the linear lower bounds of $\Phi(t)$ from Eqs.~\eqref{Eq:Phi_B1}~and~\eqref{Eq:Phi_B2}, we have
	\[
	\Phi(t,x,c)\ge
	\begin{cases}
		(\mu+k_p x)t, & 0\le t\le c,\\
		K_x +\mu(t-c), & t>c,
	\end{cases}
	\]
	where $K_x = c(\mu + k_p x)$. 
	Splitting the integral yields
	\begin{equation}
		\mathbb{E}_{x,c}[\tau_n^2] \le \int_0^c 2t \bar{F}_\varepsilon((\mu + k_p x)t) \D t + \int_c^\infty 2t \bar{F}_\varepsilon(K_x + \mu(t-c)) \D t.
	\end{equation}
	Applying change of variables $u = (\mu + k_p x)t$ to the first integral and $u = K_x + \mu(t-c)$ to the second, we obtain
	\begin{equation}
		\mathbb{E}_{x,c}[\tau_n^2] \le \frac{\int_0^{K_x} 2u \bar{F}_\varepsilon(u) \D u}{(\mu + k_p x)^2} + \frac{2c}{\mu} I(K_x) + \frac{2}{\mu^2} \int_{K_x}^\infty (u - K_x) \bar{F}_\varepsilon(u) \D u.
	\end{equation}
	Since $\mathbb{E}[\varepsilon^2] < \infty$, the first term is $\mathcal{O}(x^{-2})$. For the tail terms, the finiteness of the second moment ensures that 
	$$I(K_x) = o\left(K_x^{-1}\right), \quad\quad \int_{K_x}^\infty (u - K_x) \bar{F}_\varepsilon(u) \D u = o(K_x^{-1}).$$
	Consequently, the entire second moment term in the drift equation satisfies:
	\begin{equation}
		\frac{xp(p+1)}{2c^2} \mathbb{E}_{x,c}[\tau_n^2] = x \cdot \mathcal{O}(x^{-2}) + x \cdot o(x^{-1}) = \mathcal{O}(x^{-1}).
	\end{equation}
	Combining these, the drift satisfies:
	\begin{equation}
		\mathbb{E}_{x,c}[X_n - x] \le \xi - \frac{p \mu_{\varepsilon}}{c} \left( \frac{x}{\mu + x} \right) + \mathcal{O}(x^{-1}).
	\end{equation}
	Since $ c < 2 \gamma$, and by Eq.~\eqref{Eq:condition1}, 
	$$  -\frac{p}{c} < - \frac{p}{2\gamma}  = - \xi  - \delta, \quad\quad
	\delta = \frac{p}{2\gamma}-\xi >0.$$
	Thus, we rewrite the drift as
	$$ \mathbb{E}_{x,c}[X_n] - x < -\delta \mu_{\varepsilon} \left( \frac{x}{\mu + x} \right) +    \frac{\xi \mu}{\mu + x} + \mathcal O(x^{-1}).$$	
	Therefore, there exists \(x^*>0\) such that the drift is strictly negative for all \(x>x^*\).
\end{proof}

\begin{assumption}[Stability condition]\label{Ass:stability}
	The previous propositions provide useful guidance toward establishing sufficient conditions for stability.
	We assume that the model parameters satisfy
	\begin{equation}
		\xi < \frac{p}{2\gamma}
		\quad \text{and} \quad
		\xi < \frac{p}{C_2},
		\qquad
		C_2 = \frac{\gamma + \sqrt{\gamma^2 + \frac{3\gamma}{k_p \xi}}}{2},
		\qquad
		k_p = \frac{1-2^{\,1-p}}{p-1}.
		\label{Eq:condition}
	\end{equation}
	Equivalently, this condition can be written as
	\begin{equation}
		\xi < \min\!\left\{
		\frac{p}{2 \gamma},
		\frac{4p^2}{\gamma \left( 4p + \frac{3}{k_p} \right)} \label{Eq:condition2}
		\right\}.
	\end{equation}
	In particular, the binding constraint depends on the value of $p$:
	\begin{equation}
		\xi <
		\begin{cases}
			\displaystyle \frac{p}{2\gamma}, & \text{if } k_p p > 0.75, \\[10pt]
			\displaystyle \frac{4p^2}{\gamma \left( 4p + \frac{3}{k_p} \right)},
			& \text{if } k_p p < 0.75,
		\end{cases}
	\end{equation}
	where the critical point is approximately $p \approx 1.132$.
\end{assumption}

\section{Markov chain}\label{Sect:MC}

This section examines the process ${(X_n, c_n)}_{n \ge 0}$ as a general state-space Markov chain.
The state space of the chain is defined as
\[
\mathcal{S} = [\xi, \infty) \times [\gamma, \infty).
\]
For any current state $(x, c) \in \mathcal{S}$ and any Borel set $B \subseteq \mathcal{S}$,
the transition dynamics are characterized by the transition kernel \( P(x, c, B) \), given by
\begin{equation}
	P(x, c, B)
	= \mathbb{P} \left(
	\left( \xi + x r^{-p},\;
	\frac{\,\xi + x r^{-p}}
	{\,\xi/\gamma + (x/c)\,r^{-(p+1)}\,} \right)
	\in B
	\right),
	\label{Eq:P}
\end{equation}
where
\[
\tau = \Phi^{-1}(\varepsilon; x, c), \qquad
r = 1 + \frac{\tau}{c}, \qquad
\varepsilon \sim F_\varepsilon.
\]


\subsection{NSS and CM}

Our process $(X_n,c_n)$ is cast as a nonlinear state-space model, $\mathrm{NSS}(\Psi)$,
\[
(X_{n+1}, c_{n+1}) = \Psi (X_n, c_n, \tau_{n+1})
\]
where $\Psi$ is defined by Eqs.~\eqref{Eq:Psi},~\eqref{Eq:Psi_x} and ~\eqref{Eq:Psi_c}.
Since $\Phi(\cdot, \mathbf{X}_n)$ is strictly increasing, $\tau_{n+1}$ uniquely captures the new randomness introduced at each step, making it a valid substitute for $\varepsilon_{n+1}$ in the state-space representation.
Note that $\Psi$ is smooth in $X$, $c$ and $\tau$. 
For every $k$, the marginal law of disturbance $\tau_k$ is supported on $O_\tau=(0,\infty)\subset\mathbb R$.
We associate to $\mathrm{NSS}(\Psi)$ the controlled deterministic system
$$(X_k, c_k)=\Psi_k(x,c,u_{1:k})=\Psi\left(\Psi_{k-1}(x,c,u_{1:k-1}), u_k \right),$$
with deterministic control sequence $u_{1:k}=(u_1,\dots,u_k)$ taking values in $O_\tau$. 
We refer to this deterministic mapping as $\mathrm{CM}(\Psi)$.

We rely on standard definitions and results from \cite{meyn2009markov}, which are stated without proof; results specific to our model are proved in this paper.


\begin{definition}[Invariant set]
	Define $A_+(E)$ to be the set of all states attainable by $\mathrm{CM}(F)$ starting from any point in $E \subseteq \mathcal S$ at some time $k \geq 0$.
	If $E \subset A_+(E)$, then $E$ is called invariant.
\end{definition}

\begin{definition}[Minimal set]
	An invariant set $M$ is called a minimal set if it contains no proper closed invariant subsets. 
	That is, for any closed invariant set $E \subseteq M$, we have $E = M$.
\end{definition}

\begin{definition}[Indecomposable]
	A control system is said to be indecomposable if the state space $\mathcal S$ cannot be partitioned into two nonempty disjoint closed invariant sets.
	This implies that if a minimal set exists, it is unique.
\end{definition}



\begin{definition}[$M$-irreducibility]
	The control system is $M$-irreducible if it is indecomposable and possesses a unique minimal set $M$.
\end{definition}


\begin{definition}[Global attracting state]
	A state $x^{*} \in \mathcal S$ is called global attracting if for all $y \in \mathcal S$, we have $x^{*}\in \overline{A_{+}(y)}$.
\end{definition}

\begin{theorem}
	The nonlinear control system is $M$-irreducible if and only if a global attracting state exists.
\end{theorem}

\begin{lemma}
	The point $(\xi, \gamma)$ is a global attracting state of $\mathrm{CM}(\Psi)$.
	Thus, $\mathrm{CM}(\Psi)$ is $M$-irreducible.
\end{lemma}
\begin{proof}
	For an increasing sequence $k_n>0$ and apply the feedback control
	$$
	u_n = k_n \,c_n \quad\Longrightarrow\quad r_n = 1+\frac{u_n}{c_n} = 1 + k_n 
	$$
	and $r_n \to \infty$ as $n \to \infty$.
	The dynamics of $X_n$ is represented by
	$$
	X_n
	= \xi + \left(\prod_{i=0}^{n-1} r_i^{-p}\right)(X_0-\xi)
	+ \xi\sum_{k=0}^{n-1}\left(\prod_{i=k}^{n-1} r_i^{-p}\right),
	$$
	and choose a sufficiently fast increasing sequence $\{k_n\}$ such that
	\[
	\prod_{i=0}^{n-1} r_i^{-p}\to 0,
	\qquad 
	\sum_{k=0}^{n-1}\prod_{i=k}^{n-1} r_i^{-p} \to 0,
	\]
	then
	$ X_n \rightarrow \xi$.
	Similarly, we have $X_n / c_n \rightarrow \xi/\gamma$, it follows that $c_n = X_n / (X_n/c_n) \to \gamma.$
	Thus $(\xi,\gamma)$ is a global attracting state.
\end{proof}

\begin{definition}[Forward accessible]
	The deterministic control model $\mathrm{CM}(\Psi)$ is forward accessible if for every $(x, c) \in \mathcal S$,
	the set $A_+((x,c)) \subset \mathcal S$ has non-empty interior.	
\end{definition}

It is known that
in an $n$-dimensional model, if the generalized controllability matrix has full rank $n$ at some control sequence, then the model is forward accessible.
In our two-dimensional model, we can summarize as follows.

\begin{proposition}[Rank condition]
	For each initial condition $(x, c)$, if there exists $k \in \mathbb{Z}_{+}$ and sequence $(u_1, \cdots, u_k)$ such that the rank of the following matrix is 2:
	$$
	\frac{\partial(X_k,c_k)}{\partial(u_1,\cdots, u_k)}.
	$$
	Then the nonlinear control model $\mathrm{CM}(\Psi)$ is forward accessible.
\end{proposition}

\begin{lemma}
	$\mathrm{CM}(\Psi)$ is forward accessible.
\end{lemma}
\begin{proof}
	Using Lemma~\ref{Lemma:derivatives} and after some calculations, the determinant of two-step transition Jacobian is	
	$$
	\det\left[\frac{\partial(X_2,c_2)}{\partial(u_1,u_2)}\right]_{(u_1,u_2)=(0,0)}
	=
	\frac{p (p+1) x \xi (\gamma - c) (x+2\xi) (\gamma x + c \xi)}{c (x + \xi) (\gamma x + 2 c \xi)^2} \neq 0 \quad \text{for } c > \xi,
	$$
	see Appendix~\ref{Sect:determinant}.
	Since the map $\Psi$ is smooth, the determinant is continuous in $(u_1, u_2)$, and hence there exists $\epsilon>0$ such that
	$
	\det\big(\partial (X_2,c_2)/\partial(u_1,u_2)\big)\neq 0
	$
	for all $(u_1,u_2)\in (0,\epsilon)^2$. 
	
	If the initial state has $c=\gamma$, take any $u_1>0$.
	By continuity of the update, the first step moves the chain to a state with $c_1>\xi$, after which the two-step argument applies. 
\end{proof}

\begin{definition}[T-chain]
	Let $a$ be a sampling distribution over the non-negative integers $\mathbb{Z}_+$
	and there exists a substochastic transition kernel $T$ such that
	$$
	K_a(x, B) := \sum_{k=0}^{\infty} a(k) P^k(x, B) \geq T(x, B), \quad x \in \mathcal X, B \in \mathcal B(\mathcal X)
	$$
	where $T(\cdot, B)$ is a lower semicontinuous for any $B \in \mathcal B(\mathcal X)$.
	Then $T$ is called a continuous component of $K_a$.
	If $\{X_n\}$ is a Markov chain for which there exists a sampling distribution $a$ such that $K_a$ possesses a continuous component $T$ with $T(x, \mathcal X) > 0$ for all $x$, then $\{X_n\}$ is called a T-chain. 
\end{definition}

\begin{theorem}
	Suppose that $\mathrm{CM}(\Psi)$ is forward accessible and the disturbance process $\tau$ of $\mathrm{NSS}(\Psi)$ satisfies the density condition. 
	That is the distribution of $\tau$  has a continuous density $f_\tau$ and $O_\tau = \{x : f_\tau(x) > 0 \}$ is a open set.
	Under this condition, the following holds:
	\begin{enumerate}
		\item The $\mathrm{NSS}(\Psi)$ is a T-chain.
		\item The $\mathrm{NSS}(\Psi)$ is $\psi$-irreducible if and only if $\mathrm{CM}(\Psi)$ is $M$-irreducible.
		\item If $\mathrm{CM}(\Psi)$ is $M$-irreducible and the minimal set $M$ is aperiodic, then the $\mathrm{NSS}(\Psi)$ is a $\psi$-irreducible aperiodic T-chain.
	\end{enumerate}
\end{theorem}


\begin{lemma}[Aperiodicity of the minimal set]
	Let $\mathrm{CM}(\Psi)$ be the deterministic control model associated with our nonlinear state-space model, and let $M = \overline{A_+((\xi,\gamma))}$ be the closure of the forward reachable set from the globally attracting state $(\xi,\gamma)$. Then $M$ is aperiodic.
\end{lemma}

\begin{proof}
	Fix any neighborhood $U$ of $(\xi,\gamma)$.  
	Starting from $y=(\xi,\gamma)$, take an arbitrary control $u_1>0$; by continuity of $\Psi$, the resulting state is irrelevant.  
	Since for every $(x,c)$ we have 
	\[
	F_u(x,c)\to (\xi,\gamma)\quad\text{as }u\to\infty
	\]
	uniformly on compact sets, we may choose $u_2$ sufficiently large so that the two-step state $F_{u_2}\circ F_{u_1}(y)$ lies in $U$.  
	Similarly, for any two arbitrary controls $u_1,u_2>0$, selecting a sufficiently large $u_3$ ensures that the three-step state 
	$F_{u_3}\circ F_{u_2}\circ F_{u_1}(y)$ also lies in $U$.  
	Thus $U$ is reachable from $y$ in both 2 and 3 steps. Since $\gcd(2,3)=1$, the period of $M$ equals $1$, establishing that $M$ is aperiodic.
\end{proof}


Therefore, our model is a $\psi$-irreducible aperiodic T-chain.

\subsection{Lyapunov criterion}

\begin{theorem}[Foster--Lyapunov criterion for Joint Stability]\label{Thm:FL_Joint}
	Suppose $ \{(X_n, c_n)\}_{n \in \mathbb{N}}$ is a $\psi$-irreducible Markov chain on the state-space $\mathcal{X} = \mathbb{R}_{\ge 0} \times [\gamma, \infty)$. If there exists a measurable function $V: \mathcal{X} \to [0, \infty)$, a petite set $C \subset \mathcal{X}$, and constants $\delta > 0$ and $b < \infty$ such that
	\begin{align}
		\mathbb{E}[V(X_{n+1}, c_{n+1}) \mid X_{n} = x, c_n = c ] - V(x, c) \leq - \delta + b \mathbbm{1}_{C}(x, c), \label{Eq:Foster_Joint}
	\end{align}
	then the Markov chain is positive Harris recurrent, and there exists a unique invariant probability measure $\pi$.
\end{theorem}


\begin{theorem}
	Under the condition of Assumption~\ref{Ass:stability},
	the Markov chain $\{(X_n, c_n)\}_{n \in \mathbb{N}}$ is positive Harris recurrent.
\end{theorem}
\begin{proof}
	
	We define the Lyapunov function $V(x, c) = x + wc$ for some $w > 0$. 
	The joint drift is
	$$ \Delta V(x, c) = \mathbb{E}_{x,c}[X_{n} - x ] + w \mathbb{E}_{x,c}[c_{n} - c ].$$
	We define the compact set $C = [0, K_x] \times [\gamma, C_1]$ for some $K_x > 0$.
	By appropriately choosing $w$ and $K_x$, the drift $\Delta V(x, c)$ is negative outside of $C$.
	
	For the region $x > K_x$ and $c \le  C_2$, under the condition of Assumption~\ref{Ass:stability},
	we have
	$$ c \leq C_2  < \frac{p}{\xi}.$$
	By Proposition~\ref{Pro:ndX}, there exists $x^* >0$ such that for $x > x^*$,
	$$\mathbb{E}_{x, c}[X_n ] - x < 0.$$
	Also, by Eq.~\eqref{Eq:U_bound}, for any $K_x >0$ and $x > K_x$, we have
	\begin{align*}
		\mathbb{E}_{x,c}[\,c_n - c\,] &\leq  \dfrac{\mu_{\varepsilon}}{\mu + k_p x}
		+ \dfrac{I(c(\mu + k_p x))}{\mu}
		+ \dfrac{c\xi(\gamma - c)}{\gamma x + c \xi} \\
		&\leq   \dfrac{\mu_{\varepsilon}}{\mu + k_p K_x} +  \dfrac{I(c(\mu + k_p K_x))}{\mu}  +   \dfrac{c\xi(\gamma - c)}{\gamma K_x + c \xi}. \label{Eq:ineq_region1}
	\end{align*}
	Thus, we can  choose $K_x > x^*$ large enough such that 
	$$ \Delta V(x, c) = \mathbb{E}_{x, c}[X_n ] - x + w \left( \dfrac{\mu_{\varepsilon}}{\mu + k_p K_x} +  \dfrac{I(c(\mu + k_p K_x))}{\mu}  +   \dfrac{c\xi(\gamma - c)}{\gamma K_x + c \xi} \right) < 0 .$$
	
	Consider the regions $x < K_x$, $c > C_1$ and $x> K_x$, $c > C_2$, i.e., 
	$$[0, K_x] \times (C_1, \infty) \cup (K_x, \infty) \times (C_2, \infty).$$
	In this region, $X$ might have a positive drift but bounded by $\xi$ and by Proposition~\ref{Prop:drift_c}, $c$ has a negative drift.
	Therefore, choosing $w$ large enough, the total drift $ \Delta V$ is negative.
	
	For the compact set $C$, the one time drifts of $X$ and $c$ are bounded by $\xi$ and a constant in~Eq.~\eqref{Eq:max_cn},
	$ \Delta V $ is also bounded by some $b$.
	
\end{proof}

\section{Simulation Study}\label{Sect:simul}

To assess the long-memory properties of the proposed process, 
we conduct a simulation study with sample size $n=50{,}000$ using unit exponential distribution.
The baseline intensity is fixed at $\mu=0.4$, while the parameters $p$, $\xi$, and $\gamma$ are varied to explore different regimes.
We focus on parameter configurations located on the boundary of the stability region (Assumption~1) and in its interior, and compare the resulting dependence structures to characterize the transition from critical to stationary dynamics.
	

Persistence is quantified using the local Whittle estimator (LWE) \citep{robinson1995gaussian}.
The LWE is a semiparametric estimator of the memory parameter $d$ based on the low-frequency behavior of the spectral density, and is robust to short-run dynamics and high-frequency misspecification.


\subsection{Local Whittle estimator}

Let $\{X_t\}$ be a weakly stationary process with mean $\mu$ and autocovariance function $\gamma(k)$.
The spectral density is defined as
\begin{equation}
	f(\lambda) = \frac{1}{2\pi} \sum_{k=-\infty}^{\infty} \gamma(k) e^{-i\lambda k}, \qquad \lambda \in [-\pi,\pi].
\end{equation}
Long-memory behavior is characterized by
\begin{equation}
	f(\lambda) \sim G \lambda^{-2d} \quad \text{as } \lambda \to 0^+,
\end{equation}
where $G>0$ and $d$ is the memory parameter.
The process is stationary for $d<1/2$.

Let $x_t=\log\tau_t$ denote the log inter-arrival times.
The periodogram at Fourier frequencies $\lambda_j=2\pi j/n$ is given by
\begin{equation}
	I(\lambda_j)=\frac{1}{2\pi n}\left|\sum_{t=1}^{n}x_t e^{it\lambda_j}\right|^2,
	\qquad j=1,\dots,m,
\end{equation}
where $m<n/2$ is a bandwidth parameter.
Under a local power-law approximation of the spectral density, the LWE is defined as
\begin{equation}
	\hat d=\arg\min_{d\in\Theta} R(d),
\end{equation}
with
\begin{equation}
	R(d)=\log\!\left(\frac{1}{m}\sum_{j=1}^{m}I(\lambda_j)\lambda_j^{2d}\right)
	-\frac{2d}{m}\sum_{j=1}^{m}\log\lambda_j.
\end{equation}

The bandwidth $m$ controls a bias--variance trade-off.
Single-bandwidth estimates are reported using $m^\ast=\lfloor n^{0.6}\rfloor$, which falls within the stable plateau observed across bandwidth choices.
We also assess robustness by varying $m$ over the range $m\in[n^{0.4},\,n^{0.75}]$ and identify stable plateaus in $\hat d$.


\subsection{Simulation Results}

Table~\ref{tab:single_bw} reports single-bandwidth local Whittle estimates
for selected combinations of $p$ and $\xi$.
For each pair $(p, \xi)$, let $\gamma^*$ denote the maximum admissible $\gamma$
satisfying Assumption~\ref{Ass:stability}.
We consider three specifications: the stability boundary $\gamma = \gamma^*$,
and two interior cases $\gamma = 0.9\,\gamma^*$ and $\gamma = 0.7\,\gamma^*$.
For fixed $p$, the estimated memory parameter $\hat{d}$ increases with the
excitation parameter $\xi$, indicating that stronger excitation leads to greater
persistence.
Estimates at the stability boundary are close to the theoretical upper limit of
$0.5$, consistent with near-critical dynamics.

\begin{table}[t]
	\centering
	\caption{Single-bandwidth local Whittle estimates $\hat{d}(m^\ast)$.
	For each $(p,\xi)$, $\gamma^{\max}$ denotes the largest $\gamma$ satisfying
	the stability condition~Eq.~\eqref{Eq:condition}.
	Boundary sets $\gamma = \gamma^{*}$; Interior Cases 1 and 2
	set $\gamma = 0.9\,\gamma^{*}$ and $\gamma = 0.7\,\gamma^{*}$, respectively.}
	\label{tab:single_bw}
	\renewcommand{\arraystretch}{1.3}
	\begin{tabular}{c c c c | c c c | c c c}
		\hline
		& \multicolumn{3}{c|}{Stability Boundary} & \multicolumn{3}{c|}{Interior Case 1 (0.9$\gamma^{*}$)} & \multicolumn{3}{c}{Interior Case 2 (0.7$\gamma^{*}$)} \\
		\cline{2-10}
		$p$ & $\xi$ & $\gamma$ & $\hat{d}(m^\ast)$ & $\xi$ & $\gamma$ & $\hat{d}(m^\ast)$ & $\xi$ & $\gamma$ & $\hat{d}(m^\ast)$ \\
		\hline
		1.001 & 0.1 & 4.8094 & 0.2195 & 0.1 & 4.3285 & 0.2875 & 0.1 & 3.3666 & 0.2330 \\
		& 0.5 & 0.9618 & 0.4648 & 0.5 & 0.8657 & 0.3971 & 0.5 & 0.6733 & 0.2007 \\
		& 1.0 & 0.4809 & 0.4958 & 1.0 & 0.4328 & 0.4327 & 1.0 & 0.3367 & 0.1341 \\
		\hline
		1.1   & 0.1 & 5.4505 & 0.2711 & 0.1 & 4.9055 & 0.2996 & 0.1 & 3.8150 & 0.2598 \\
		& 0.5 & 1.0901 & 0.4540 & 0.5 & 0.9811 & 0.4133 & 0.5 & 0.7631 & 0.1370 \\
		& 1.0 & 0.5451 & 0.4765 & 1.0 & 0.4905 & 0.4131 & 1.0 & 0.3815 & 0.1053 \\
		\hline
		1.5   & 0.1 & 7.5000 & 0.3024 & 0.1 & 6.7500 & 0.2964 & 0.1 & 5.2500 & 0.1501 \\
		& 0.5 & 1.5000 & 0.3899 & 0.5 & 1.3500 & 0.2698 & 0.5 & 1.0500 & 0.0985 \\
		& 1.0 & 0.7500 & 0.3242 & 1.0 & 0.6750 & 0.1979 & 1.0 & 0.5250 & 0.0601 \\
		\hline
		2.0   & 0.1 & 10.000 & 0.2933 & 0.1 & 9.0000 & 0.2152 & 0.1 & 7.0000 & 0.0921 \\
		& 0.5 & 2.0000 & 0.2457 & 0.5 & 1.8000 & 0.1561 & 0.5 & 1.4000 & 0.0335 \\
		& 1.0 & 1.0000 & 0.2096 & 1.0 & 0.9000 & 0.1152 & 1.0 & 0.7000 & 0.0043 \\
		\hline
	\end{tabular}
\end{table}


The stability of the estimator is visually confirmed in Figure \ref{fig:lwe_stability}, where $\hat{d}$ remains remarkably consistent across a wide range of bandwidths. 
This lack of sensitivity to $m$ validates the robustness of our long-memory inference and suggests that the observed power-law behavior is a structural feature of the simulated process rather than an artifact of bandwidth selection.

\begin{figure}[t]
	\centering
	\includegraphics[width=0.75\textwidth]{local_whittle_p_1.1_xi_0.5_b.pdf}
	\caption{Stability plot of the Local Whittle estimates $\hat{d}(m)$ as a function of the bandwidth $m$ (Parameters: $p=1.1, \xi=0.5, \gamma = 1.0901$).}
	\label{fig:lwe_stability}
\end{figure}


\section{Conclusion}~\label{Sect:conc}

We proposed a self-exciting point process with power-law intensity dynamics that admits a finite-dimensional Markovian representation.
The model preserves the local update structure of power-law Hawkes processes while replacing global history dependence with a nonlinear Markov chain governing the latent intensity states.
Within a general state-space framework, we established irreducibility, aperiodicity, and positive Harris recurrence under explicit stability conditions, ensuring long-run stability of the process.
Simulation results based on the local Whittle estimator show that the proposed model exhibits long-memory behavior near the boundary of the stability region.

\section*{Funding}
This work was supported by the National Research Foundation of Korea (NRF) grant funded by the Korea government (MSIT) (No. RS-2026-25469087).




\bibliography{Bib}
\bibliographystyle{chicago}


\appendix

\section{Derivation of determinant}\label{Sect:determinant}

We use the following notations:
\begin{align}
	X_n = X_{n-1} r_{n}^{-p} + \xi, \qquad D_n = \frac{X_{n-1}}{c_{n-1}}r_n^{-(p+1)}+\frac{\xi}{\gamma}, \qquad r_n = 1 + \frac{u_n}{c_{n-1}}, \qquad c_n = \frac{X_n}{D_n}
\end{align}
with $X_0 = x$ and $c_0 = c$.
Moreover, for any function $g$ of $u_i$, the notation $\left. g(u_i) \right|_{\mathbf{0}}$ indicates that the function is evaluated at $u_i = 0$ for all $i$.
Since $\left. r_n \right|_{\mathbf{0}}=1$, we have
\begin{equation}
	\left. X_1 \right|_{\mathbf{0}} = x + \xi,  \qquad
	\left. D_1 \right|_{\mathbf{0}} = \frac{\gamma x + c \xi}{\gamma c}, \qquad \left. c_1 \right|_{\mathbf{0}} = \frac{\gamma c(x + \xi)}{\gamma x + c \xi}, \label{Eq:1_0}
\end{equation}
and
\begin{equation}
\left. X_2 \right|_{\mathbf{0}} = x + 2 \xi,\qquad \left. D_2 \right|_{\mathbf{0}} = \frac{\gamma x + 2 c \xi}{\gamma c}, \quad \left. c_2 \right|_{\mathbf{0}} = \frac{\gamma c(x + 2\xi)}{\gamma x + 2 c \xi}. 
\label{Eq:2_0}
\end{equation}

\begin{lemma}\label{Lemma:derivatives}
The evaluated derivatives of of $X_2$ and $c_2$ with respect to $u_i$ at $u_1 = u_2 =0$ are
\begin{align}
&\left.\frac{\partial X_2}{\partial u_1}\right|_{\mathbf{0}} = -\frac{px}{c}, \quad\quad\quad\quad \left.\frac{\partial X_2}{\partial u_2}\right|_{\mathbf{0}} =  -\frac{p(\gamma x+c\xi)}{\gamma c}, \label{Eq:eval_1}\\
&\left.\frac{\partial c_2}{\partial u_1}\right|_{\mathbf{0}}
= \frac{\gamma x\big(\gamma x + 2p\xi(\gamma-c) + 2\gamma\xi\big)}{(2c\xi+\gamma x)^2} , \quad\quad\quad\quad
\left.\frac{\partial c_2}{\partial u_2}\right|_{\mathbf{0}}
= (\gamma x + c\xi) \frac{B}{M}	 \label{Eq:eval_2}
\end{align}
where
\begin{align*}
&B = \gamma x^2 + (p+2)\gamma \xi x + (1-p)c\xi x + 2c\xi^2, \\
&M = (x+\xi) (\gamma x + 2 c\xi)^2.
\end{align*}
\end{lemma}
\begin{proof}
	We begin with Eq.~\eqref{Eq:eval_1}.
	First, note that
	$$ \frac{\partial X_2}{\partial c_1} = p X_1 r_2^{-p-1} u_2 c_1^{-2}.$$
	Hence,
	\begin{equation}
	\left.\frac{\partial X_2}{\partial u_1}\right|_{\mathbf{0}}
	=
	\left.\frac{\partial X_2}{\partial X_1}\right|_{\mathbf{0}}\left.\frac{\partial X_1}{\partial u_1}\right|_{\mathbf{0}}
	+
	\left.\frac{\partial X_2}{\partial c_1}\right|_{\mathbf{0}}\left.\frac{\partial c_1}{\partial u_1}\right|_{\mathbf{0}}
	=
	1\cdot\left(-\frac{p x}{c}\right)+0\cdot(\cdots)
	=
	-\frac{p x}{c}, \label{Eq:X2_u1_0}
	\end{equation}
	Similarly,
	\begin{equation}
	\frac{\partial X_2}{\partial u_2} = 
	X_1\cdot(-p) r_2^{-(p+1)}\cdot\frac{1}{c_1}
	\quad\Rightarrow\quad
	\left.\frac{\partial X_2}{\partial u_2}\right|_{\mathbf{0}}
	= -\frac{p X_1}{c_1}\bigg|_{\mathbf{0}}
	= -\frac{p(\gamma x+c\xi)}{\gamma c}. \label{Eq:X2_u2_0}
	\end{equation}

	For the first expression in Eq.~\eqref{Eq:eval_2}, observe that
	\begin{equation}
	\left.\frac{\partial c_2}{\partial u_1}\right|_{\mathbf{\mathbf{0}}}
	= \left.\frac{\partial c_2}{\partial X_1}\right|_{\mathbf{0}}\left.\frac{\partial X_1}{\partial u_1}\right|_{\mathbf{0}} + 
	\left.\frac{\partial c_2}{\partial c_1}\right|_{\mathbf{0}}\left.\frac{\partial c_1}{\partial u_1}\right|_{\mathbf{0}}. \label{Eq:dc2_du1}
	\end{equation}
	From Eqs.~\eqref{Eq:1_0}, \eqref{Eq:2_0}, and
	\begin{align*}
	&\frac{\partial {X_2}}{\partial X_1} = r_2^{-p} \quad\Rightarrow\quad  \left. \frac{\partial {X_2}}{\partial X_1} \right|_{\mathbf{0}} = 1 \\ 
	& \frac{\partial{D_2}}{\partial X_1} =  - D_1 (p+1) r_2^{-(p+2)}\frac{\partial r_2}{\partial X_1} =  (p+1) r_2^{-(p+2)}\frac{u_2 D_1^2}{X_1^2} \quad\Rightarrow\quad  \left. \frac{\partial {D_2}}{\partial X_1} \right|_{\mathbf{0}} = 0,
	\end{align*}
	we obtain
	\begin{align}
		\left.\frac{\partial c_2}{\partial X_1}\right|_{\mathbf{0}} = \left. \frac{(\partial_{X_1}X_2) D_2 - X_2 (\partial_{X_1}D_2)}{D_2^2} \right|_{\mathbf{0}} = \left. \frac{1}{D_2} \right|_{\mathbf{0}} =  \frac{\gamma c}{\gamma x + 2c\xi }. \label{Eq:dc2_dX1}
	\end{align} 
	Also,
	\begin{equation}
	\left. \frac{\partial X_1}{\partial u_1} \right|_{\mathbf{{0}}} = \left. - x p r_{1}^{-p-1}\frac{1}{c} \right|_{\mathbf{{0}}}= -\frac{px}{c}. \label{Eq:dX_1_du_1_0}
	\end{equation}
	
	Next, since
	\begin{align*}
		\frac{\partial c_2}{\partial c_1} = \frac{\partial}{\partial c_1} \left(\frac{X_2}{D_2}\right) = \frac{(\partial_{c_1}X_2) D_2 - X_2 (\partial_{c_1}D_2)}{D_2^2} ,
	\end{align*}
	and
	$$ \left. (\partial_{c_1}X_2) \right|_{\mathbf{0}} = \left. pX_1 r_2^{p-1} u_2 \frac{1}{c_1^{2}} \right|_{\mathbf{0}} = 0, \qquad  \left. (\partial_{c_1}D_2) \right|_{\mathbf{0}} = \left. - \frac{X_1}{c_1^2} \right|_{0} = - \frac{(\gamma x + c \xi)^2}{(c\gamma)^2 (x + \xi)},
	$$
	by Eq.~\eqref{Eq:2_0},
	we have
	\begin{equation}
	\left.\frac{\partial c_2}{\partial c_1}\right|_{\mathbf{0}}
	= \frac{(x+2\xi)(\gamma x + c\xi)^2}{(x+\xi)(\gamma x + 2c\xi)^2 }. \label{Eq:dc2_dc1}
	\end{equation}
	From Eqs.~\eqref{Eq:1_0} and \eqref{Eq:dX_1_du_1_0} and using
	$$\left. \partial_{u_1} D_1 \right|_{\mathbf{0}} = - \frac{(p+1)x}{c^2},$$
	we have
	$$
	\left. D_1 (\partial_{u_1} X_1) - X_1 (\partial_{u_1} D_1) \right|_{\mathbf{0}}
	=  -\frac{px}{c} \frac{x \gamma + \xi c}{c \gamma} + \frac{(p+1)x}{c^2}(x + \xi )
	$$
	and
	\begin{equation}
	\left.\frac{\partial c_1}{\partial u_1}\right|_{\mathbf{0}}
	=\left. \frac{D_1 (\partial_{u_1} X_1) - X_1 (\partial_{u_1} D_1) }{D_1^2} \right|_{0}=  \frac{x\gamma\big(\gamma x + (p+1)\gamma\xi - p c\xi\big)}{(\gamma x + c \xi)^2}. \label{Eq:dc1_du1}
	\end{equation}
	Substituting Eqs.~\eqref{Eq:dc2_dX1},\eqref{Eq:dX_1_du_1_0},\eqref{Eq:dc2_dc1} and \eqref{Eq:dc1_du1} into Eq.~\eqref{Eq:dc2_du1} yields the desired expression for $\partial c_2 / \partial u_1$.

	For the second equality in Eq.~\eqref{Eq:eval_2}, note that
	$$
	\frac{\partial c_2}{\partial u_2}
	=\frac{(\partial_{u_2}X_2) D_2 - X_2 (\partial_{u_2}D_2)}{D_2^2}.
	$$
	From Eq.~\eqref{Eq:X2_u2_0} and
	$$
	\partial_{u_2} D_2 = \frac{X_1}{c_1} (-(p+1)) r_2^{-(p+2)}\frac{1}{c_1}
	\quad\Rightarrow\quad
	\left.\partial_{u_2}D_2\right|_{\mathbf{0}} = \left. -\frac{(p+1) X_1}{c_1^2} \right|_{\mathbf{0}},
	$$
	we find
	\begin{align*}
	 \left. (\partial_{u_2}X_2)D_2 - X_2(\partial_{u_2}D_2) \right|_{\mathbf{0}}
	&= \left. \left(-\frac{p X_1}{c_1}\right)D_2 - (X_1+\xi)\left(-\frac{(p+1) X_1}{c_1^2}\right) \right|_{\mathbf{0}} \\
	&= \left. \frac{X_1}{c_1}\left( (p+1) \frac{X_1 + \xi}{c_1} - p D_2 \right) \right|_{\mathbf{0}} \\
	&=  \frac{\gamma x + c\xi}{c \gamma } 
	\left( (p+1)\frac{(x+2\xi)(\gamma x + c\xi)}{c\gamma(x+\xi)} -p\frac{\gamma x+2c\xi}{c \gamma}\right).
	\end{align*}
	Thus,
	$$
	\left.\frac{\partial c_2}{\partial u_2}\right|_{\mathbf{0}}
	= \frac{(\gamma x + c\xi)}{(x+\xi)(\gamma x + 2c\xi)^2}
	\Big( \gamma x^2 + ((p+2)\gamma + (1-p)c ) \xi x + 2c\xi^2 \Big),
	$$
	which completes the proof.
\end{proof}


Finally, for the determinant,
$$
\det J\Big|_{\mathbf{0}}
= \left.\frac{\partial X_2}{\partial u_1}\right|_{\mathbf{0}}
\left.\frac{\partial c_2}{\partial u_2}\right|_{\mathbf{0}} - 
\left.\frac{\partial X_2}{\partial u_2}\right|_{\mathbf{0}}
\left.\frac{\partial c_2}{\partial u_1}\right|_{\mathbf{0}}
$$
and
$$\det J\Big|_{\mathbf{0}} = \frac{p (p+1) x \xi (\gamma - c) (x+2\xi) (\gamma x + c \xi)}{c (x + \xi) (\gamma x + 2 c \xi)^2}.$$



\end{document}
